% ==========================================================
% titles/title-05.tex -- Title V: Executive Function
% ==========================================================

\ConstTitle{Executive Function}{title:executive}

\Article{Principles of Executive Design}{art:executive-design}

There is no President, Prime Minister, Chancellor, or any single executive office with general authority over the Commonwealth. Executive function is distributed across domain-specific bodies with no cross-domain executive power. The specific combination of cross-domain authority and discretionary power that enables authoritarian consolidation does not exist in this Constitution.

\Article{Implementation Councils}{art:implementation-councils}

\ArtSection{Structure}{art:implementation-councils:s-structure}
Each major governance domain -- including healthcare, education, food systems, infrastructure, environment, housing, justice, and economic coordination -- has an \gls{gls:implementation-councils} responsible for translating legislative decisions into operational reality. The domain list is illustrative, not exhaustive.

\ArtSection{Leadership}{art:implementation-councils:s-leadership}
Each Council is led by a Director selected by sortition from a pool of persons with demonstrated domain competence, assessed through the \gls{gls:qualification-commons} under Article \ref{art:sortition-design}, \ref{art:sortition-design:s-qualification-commons}. Directors serve four-year non-renewable terms. No Director holds authority outside their domain.

\ArtSection{Mandate Constraint and Logging}{art:implementation-councils:s-mandate-constraint-logging}
\gls{gls:implementation-councils}s implement what the Legislature has decided; they do not make policy. All implementation decisions at or above the significance thresholds defined in Article \ref{art:transparency-engine}, \ref{art:transparency-engine:s-significance-thresholds-logging} are automatically logged in real time by the \gls{gls:transparency-engine}.

\ArtSection{Cross-Domain Coordination}{art:implementation-councils:s-cross-domain-coordination}
A \gls{gls:coordination-assembly} composed of one rotating representative from each \gls{gls:implementation-councils} meets publicly on a monthly basis to manage cross-domain issues. Each Council's representative rotates annually by internal lottery. Extraordinary sessions may be convened by any three member councils on fourteen days' notice. The Assembly may only flag conflicts and propose solutions to the Legislature; it makes no binding decisions.

\ArtSection{Creation and Dissolution}{art:implementation-councils:s-creation-dissolution}
The \gls{gls:sortition-legislature} may create new \gls{gls:implementation-councils}s by supermajority vote when new governance domains emerge. Dissolution requires supermajority vote and \gls{gls:constitutional-court} review confirming that no Material Sufficiency Floor function is left without an implementing body.

\Article{Emergency Protocols}{art:emergency-protocols}

\ArtSection{Pre-Legislation Requirement}{art:emergency-protocols:s-pre-legislation-requirement}
For each predictable category of emergency -- including pandemic, major natural disaster, infrastructure failure, external military attack, and economic crisis -- the \gls{gls:sortition-legislature} pre-approves a detailed response protocol during normal times, free from crisis pressure. For each predictable cross-domain emergency category, the Legislature simultaneously designates a lead \gls{gls:implementation-councils} and supporting Councils with defined coordination roles. Lead and supporting designations are published alongside their protocols through the \gls{gls:transparency-engine}.

\ArtSection{Automatic Activation}{art:emergency-protocols:s-automatic-activation}
Emergency protocols activate automatically upon objective triggering conditions verified by the \gls{gls:emergency-threshold-verification-body}. No individual may declare an emergency to unlock expanded powers. The power to declare an emergency is not vested in any person.

The \gls{gls:emergency-threshold-verification-body} is a standard body of sixteen domain experts serving four-year terms, operating in continuous session. It monitors objective metrics for each emergency protocol category and publicly reports when thresholds are crossed. It makes no policy decisions. It cannot be instructed to withhold or delay a threshold determination.

\ArtSection{Hard Limits on Emergency Powers}{art:emergency-protocols:s-hard-limits-emergency}
No emergency protocol may, under any circumstances:
\begin{itemize}
  \item Suspend or reduce the Material Sufficiency Floor (Article \ref{art:matsuff})
  \item Override the Anti-Domination Principle (Article \ref{art:antidom})
  \item Suppress the \gls{gls:transparency-engine} (Article \ref{art:transparency-engine}), the \gls{gls:whistleblower-infrastructure} (Article \ref{art:transparency-engine}, \ref{art:transparency-engine:s-whistleblower-infrastructure}), or the \gls{gls:civic-initiative-infrastructure} (Article \ref{art:civic-initiative-infrastructure}), including its Emergency Protocol Post-hoc Review function
  \item Extend beyond its pre-defined time limit without legislative reauthorization
  \item Grant any person authority beyond the specific domain of the emergency
\end{itemize}

\ArtSection{Novel Emergencies: The Three-Part Gap Protocol}{art:emergency-protocols:s-novel-emergencies-three}
\textbf{Gap Coverage:} The \gls{gls:emergency-threshold-verification-body} determines which domain faces the largest immediate threat to Material Sufficiency Floor delivery, using a published methodology developed during normal times in consultation with all \gls{gls:implementation-councils}s, reviewed every five years alongside the climate mandate under Article \ref{art:ecological-limits}. The Council whose domain faces the largest immediate threat activates the most conservative available adjacent protocol as lead Council. Supporting Councils with adjacent jurisdiction activate complementary protocols without contradiction. If the Material Sufficiency Floor impact determination does not produce a clear lead Council within six hours of emergency threshold confirmation, the \gls{gls:residual-coordination-council} assumes narrow binding coordination authority strictly limited to:
\begin{itemize}
  \item Designating a lead Council from among those with active jurisdiction
  \item Ensuring adjacent protocol activation without contradiction
  \item Convening the seventy-two-hour legislative session
\end{itemize}
This authority expires automatically and without residue when the Legislature convenes. The \gls{gls:residual-coordination-council} may not direct Council operations, override Council decisions, or act beyond these three functions under any circumstances.

\textbf{Emergency Session:} The \gls{gls:sortition-legislature} convenes within seventy-two hours with a quorum of at least one third of members. It may authorize specific time-limited actions beyond the adjacent protocol's scope. These authorizations expire after thirty days unless actively reauthorized.

\textbf{Mandatory Review:} The full Legislature reconvenes within thirty days to review all actions taken under emergency authorization. No emergency authorization may be extended without this active legislative reauthorization.

The seventy-two-hour gap is a genuine vulnerability in this Constitution. Its design minimizes but cannot eliminate the risk of insufficient response in its early hours. The hard limits of \ref{art:emergency-protocols:s-hard-limits-emergency} apply in full even during this gap, including during any period of \gls{gls:residual-coordination-council} fallback authority.

\Article{The Residual Coordination Council}{art:residual-coordination}

\ArtSection{Composition}{art:residual-coordination:s-composition}
A \gls{gls:residual-coordination-council} of five persons, selected by sortition from former \gls{gls:implementation-councils} Directors, former \gls{gls:foreign-relations-council} members, and former \gls{gls:constitutional-court} justices who have completed their terms in good standing and are not currently serving in any governance role, serves two-year non-renewable terms. Decisions require unanimity of all five members.

If the eligible pool cannot produce five willing and unconflicted members, the \gls{gls:electoral-commons} conducts emergency sortition from all former governance body members who have completed terms in good standing. If this pool is also insufficient, the \gls{gls:constitutional-court} designates temporary coordinators from its membership solely to perform the legislative convening function of Article \ref{art:emergency-protocols}, \ref{art:emergency-protocols:s-novel-emergencies-three}. This designation expires automatically when the Legislature convenes and confers no other \gls{gls:residual-coordination-council} authority.

\ArtSection{Permitted Functions}{art:residual-coordination:s-permitted-functions}
The Council's authority is strictly limited to:
\begin{itemize}
  \item Flagging coordination failures between domains to the Legislature
  \item Receiving and responding to \gls{gls:transparency-engine} anomaly reports that identify coordination failures crossing domain boundaries, and referring findings to the Legislature with recommended remediation
  \item Representing the Commonwealth in multilateral bodies where a single voice is required, in coordination with the \gls{gls:foreign-relations-council} (Article \ref{art:foreign-relations})
  \item Convening the rapid-response legislative session in novel emergencies as defined in Article \ref{art:emergency-protocols}
\end{itemize}

In multilateral settings, the \gls{gls:foreign-relations-council} holds primacy over all matters of foreign policy substance. Substance encompasses any position, commitment, or statement that could create legal obligation, establish Commonwealth policy, or be interpreted by other parties as binding. The \gls{gls:residual-coordination-council} holds primacy over logistics, encompassing representation arrangements, procedural votes, scheduling, and forum participation decisions that do not themselves create policy commitments. When either body disputes the categorization of a specific matter, the \gls{gls:constitutional-court} (Article \ref{art:const-court}) makes a binding determination within seventy-two hours. The \gls{gls:residual-coordination-council} may not articulate any substantive position in any multilateral body without prior \gls{gls:foreign-relations-council} agreement.

In cross-domain novel emergencies where neither pre-designation nor the Material Sufficiency Floor impact determination produces a clear lead \gls{gls:implementation-councils} within six hours of threshold confirmation, the \gls{gls:residual-coordination-council} assumes the narrow binding coordination authority defined in Article \ref{art:emergency-protocols}, \ref{art:emergency-protocols:s-novel-emergencies-three}. This authority is strictly time-limited, automatically expires when the Legislature convenes, and does not expand or modify any other function of the \gls{gls:residual-coordination-council}.

It may not make any binding policy decision, direct any \gls{gls:implementation-councils}s, declare any emergency, or issue any executive order.